\chapter*{Abstract}
This thesis investigates the applicability of mechanical finite element (\acf{FEM}) simulations to predict the tensile stress and mass of chemical samples deposited on perforated nanomechanical membranes, as used for nanoelectromechanical systems-based infrared spectroscopy. The membranes are modeled as prestressed elastic structures, and their eigenfrequencies are computed and compared with experimental measurements for different deposited materials and masses. 

A calibration of the simulation framework was performed using unsampled membranes from various manufacturers and with different prestresses and sizes, demonstrating agreement within the error bounds of the measurement ($<1\%$ with experimental data across a broad parameter range. For sampled membranes, systematic measurements of eigenfrequencies were carried out with polystyrene (PS) and polypropylene (PP) deposits spanning masses from \SI{625}{\pico\gram} to \SI{40}{\nano\gram}. The results reveal a strong material dependence: PS samples induce measurable frequency shifts consistent with mass loading and negligible prestress, whereas PP samples show almost no shifts, which can be explained only by introducing an internal prestress in the deposited material. 

Quantitative agreement between simulation and experiment was achieved for most cases; however, the predictive power of purely mechanical modeling is not testable for small sample masses below \SI{10}{\nano\gram}. In this regime, the frequency shifts are smaller than experimental uncertainties. Nonetheless, the developed simulation framework enables reliable estimation of membrane prestress and offers a basis for future extensions, such as constructing simulation databases to infer mass distributions from frequency response data. 

A challenge in modeling these thin membranes are different length scales in the thickness and width of the membrane. Previous modeling attempts in Comsol have shown problems with high aspect ratios of the finite elements. This was solved by using a 2D simulation. Overall, this work demonstrates both the potential and the limitations of \acf{FEM}-based modeling of nanomechanical membranes with deposited samples, providing insights for further development of quantitative analysis methods in nanomechanical spectroscopy.

\newpage
\mbox{}
\thispagestyle{empty}

\chapter*{Zusammenfassung}
Diese Arbeit untersucht die Anwendbarkeit mechanischer Finite-Elemente-Simulationen (\acf{FEM}) zur Vorhersage der Zugspannung der Membran und Masse chemischer Proben, die auf perforierten nanomechanischen Membranen aufgebracht werden, wie sie in nanoelektromechanikbasierten Infrarotspektroskopiesystemen verwendet werden. Die Membranen werden als vorgespannte elastische Strukturen modelliert, und ihre Eigenfrequenzen werden berechnet und mit experimentellen Messungen für unterschiedliche aufgebrachte Materialien und Massen verglichen.

Eine Kalibrierung des Simulationsframeworks wurde mithilfe von Membranen ohne aufgebrachte Proben verschiedener Hersteller sowie unterschiedlicher Vorspannungen und Größen durchgeführt. Dabei zeigte sich eine Übereinstimmung innerhalb der Messfehlergrenzen (Abweichung kleiner als $1\%$) über einen breiten Parameterbereich. Für mit Proben belegte Membranen wurden systematische Eigenfrequenzmessungen mit Polystyrol- (PS) und Polypropylen- (PP) Ablagerungen im Massenbereich von \SI{625}{\pico\gram} bis \SI{40}{\nano\gram} durchgeführt. Die Ergebnisse zeigen eine starke Materialabhängigkeit: PS-Proben verursachen messbare Frequenzverschiebungen, die mit Massebelastung und vernachlässigbarer Vorspannung der Proben konsistent sind, während PP-Proben nahezu keine Verschiebungen zeigen. Dieses Verhalten lässt sich nur durch die Einführung einer inneren Vorspannung im aufgebrachten Material erklären.

In den meisten Fällen konnte eine quantitative Übereinstimmung zwischen Simulation und Experiment erreicht werden; die Vorhersagekraft rein mechanischer Modelle erwies sich jedoch für kleine Probenmassen unterhalb von \SI{10}{\nano\gram} als begrenzt. In diesem Bereich sind die Frequenzverschiebungen kleiner als die experimentellen Unsicherheiten. Dennoch ermöglicht das entwickelte Simulationsframework eine zuverlässige Abschätzung der Membranvorspannung und bietet eine Grundlage für zukünftige Erweiterungen, etwa zur Erstellung von Simulationsdatenbanken, um Massenverteilungen aus Frequenzantwortdaten abzuleiten.

Eine Herausforderung bei der Modellierung dieser dünnen Membranen stellen die unterschiedlichen Längenskalen in Dicke und Breite der Membran dar. Frühere Modellierungsansätze in Comsol zeigten Probleme aufgrund hoher Seitenverhältnisse der finiten Elemente. Dieses Problem wurde durch die Verwendung einer zweidimensionalen Simulation gelöst. Insgesamt zeigt diese Arbeit sowohl das Potenzial als auch die Grenzen \acf{FEM}-basierter Modellierung nanomechanischer Membranen mit abgeschiedenen Proben auf und liefert Erkenntnisse für die Weiterentwicklung quantitativer Analysemethoden in der nanomechanischen Spektroskopie. 


\newpage
\mbox{}
\thispagestyle{empty}


\chapter*{Acknowledgement}

I would like to express my sincere gratitude to my supervisors for their guidance, support, and valuable contributions throughout the course of this thesis. I am especially grateful to Markus Faustmann, whose expertise and advice were of great help in developing the mathematical modelling aspects of this work. I would also like to thank Niklas Luhmann for sharing his knowledge and experience in experimental scientific practice, and for teaching me a great deal about rigorous and thoughtful scientific work. I would also like to thank Jelena Timarac-Popović for her generous assistance in collecting the experimental data and for preparing the samples used in this work.

I am also deeply thankful to my fellow students for the enriching experiences we shared throughout this master’s programme, as well as for their friendship, support, and the many valuable discussions along the way.

Finally, I would like to thank everyone involved in creating and maintaining the Computational Science and Engineering master’s programme at TU Wien. I have thoroughly enjoyed my time in this programme, and it has greatly expanded both my knowledge and my abilities as a scientist and engineer.